% GENERATED FILE. Edit canonical lessons, then run npm run generate:books.
% canonical-source-hash: fnv1a64:f34b4f306aaa4a10
% canonical-lessons: SA-C384-udyamah, SA-C384-dhrtih, SA-C384-yathasakti, SA-C384-anutisthati, SA-C384-protsahayati

\chapter{Effort}
\label{ch:sa-effort}
\hlchaptermodality{384}

\begin{chapteropening}
\textbf{By the end of this chapter:} \emph{I can say what effort it will take to succeed.}

Complete the last lesson of chapter 384: I can say what effort it will take to succeed.
\end{chapteropening}

\section[\texorpdfstring{\sk{उद्यमः} (udyamaḥ)}{उद्यमः (udyamaḥ)}]{\sk{उद्यमः} (udyamaḥ) — perseverance; exertion}

\label{lesson:SA-C384-udyamah}



\begin{quote}
Before the new one: say the Sanskrit for a heartfelt wish, then the Sanskrit for more and more.
\end{quote}

\subsection*{You'll want to know: \sk{उद्यमः}}
\textbf{\sk{उद्यमः}} — \emph{udyamaḥ} — \textquotedblleft{}perseverance; exertion\textquotedblright{}.

\begin{quote}
\sk{यदि} \sk{उद्यमः} \sk{अस्ति}, \sk{तर्हि} \sk{सफलता} \sk{अपि} \sk{भविष्यति।}
\end{quote}

\begin{quote}
\emph{yadi udyamaḥ asti, tarhi saphalatā api bhaviṣyati.}
\end{quote}

\begin{quote}
If there is perseverance, success will follow.
\end{quote}

A masculine noun.

\begin{cousinweb}
\textbf{\sk{उद्यमः}} (\emph{udyamaḥ}) is built on \emph{ud}, \textquotedblleft{}up\textquotedblright{}, and the root √yam, \textquotedblleft{}to hold\textquotedblright{}; MW gives \emph{ud-yama} \textquotedblleft{}undertaking\textquotedblright{} and \textquotedblleft{}exertion, strenuous and continued effort, perseverance\textquotedblright{} (MW).
\end{cousinweb}

\begin{grammarlens}[title={What you've built}]
One more word for saying what happens if.
\end{grammarlens}

\subsection*{Your turn}
\begin{itemize}
\raggedright
  \item \emph{Say it:} \emph{udyamaḥ}
  \item \emph{Say it:} \emph{udyamaḥ}, once more
  \item \emph{Say it:} \emph{yadi udyamaḥ asti, tarhi saphalatā api bhaviṣyati.} — If there is perseverance, success will follow.
  \item \emph{Recall:} say the Sanskrit for a heartfelt wish, then the Sanskrit for more and more, then say \emph{udyamaḥ} again
\end{itemize}

\begin{tcolorbox}[breakable,colback=teal!4,colframe=teal!35!black,title={Before you move on}]
What does \sk{उद्यमः} mean? (\textquotedblleft{}Perseverance; exertion\textquotedblright{}.) Say it once more.
\end{tcolorbox}
\section[\texorpdfstring{\sk{धृतिः} (dhṛtiḥ)}{धृतिः (dhṛtiḥ)}]{\sk{धृतिः} (dhṛtiḥ) — constancy; resolve}

\label{lesson:SA-C384-dhrtih}



\begin{quote}
Before the new one: say the Sanskrit for more and more, then the Sanskrit for perseverance.
\end{quote}

\subsection*{You'll want to know: \sk{धृतिः}}
\textbf{\sk{धृतिः}} — \emph{dhṛtiḥ} — \textquotedblleft{}constancy; resolve\textquotedblright{}.

\begin{quote}
\sk{यदि} \sk{धृतिः} \sk{अस्ति}, \sk{तर्हि} \sk{भयं} \sk{न} \sk{भवति।}
\end{quote}

\begin{quote}
\emph{yadi dhṛtiḥ asti, tarhi bhayaṃ na bhavati.}
\end{quote}

\begin{quote}
If there is constancy, there is no fear.
\end{quote}

A feminine noun.

\begin{cousinweb}
\textbf{\sk{धृतिः}} (\emph{dhṛtiḥ}) comes from the root √dhṛ, \textquotedblleft{}to hold\textquotedblright{}; MW gives \emph{dhṛti} \textquotedblleft{}firmness, constancy, resolution, will\textquotedblright{} (MW).
\end{cousinweb}

\begin{grammarlens}[title={What you've built}]
One more word for saying what happens if.
\end{grammarlens}

\subsection*{Your turn}
\begin{itemize}
\raggedright
  \item \emph{Say it:} \emph{dhṛtiḥ}
  \item \emph{Say it:} \emph{dhṛtiḥ}, once more
  \item \emph{Say it:} \emph{yadi dhṛtiḥ asti, tarhi bhayaṃ na bhavati.} — If there is constancy, there is no fear.
  \item \emph{Recall:} say the Sanskrit for more and more, then the Sanskrit for perseverance, then say \emph{dhṛtiḥ} again
\end{itemize}

\begin{tcolorbox}[breakable,colback=teal!4,colframe=teal!35!black,title={Before you move on}]
What does \sk{धृतिः} mean? (\textquotedblleft{}Constancy; resolve\textquotedblright{}.) Say it once more.
\end{tcolorbox}
\section[\texorpdfstring{\sk{यथाशक्ति} (yathāśakti)}{यथाशक्ति (yathāśakti)}]{\sk{यथाशक्ति} (yathāśakti) — as far as one can}

\label{lesson:SA-C384-yathasakti}



\begin{quote}
Before the new one: say the Sanskrit for perseverance, then the Sanskrit for constancy.
\end{quote}

\subsection*{You'll want to know: \sk{यथाशक्ति}}
\textbf{\sk{यथाशक्ति}} — \emph{yathāśakti} — \textquotedblleft{}as far as one can\textquotedblright{}.

\begin{quote}
\sk{यदि} \sk{भवान्} \sk{यथाशक्ति} \sk{यतते}, \sk{तर्हि} \sk{अलम्।}
\end{quote}

\begin{quote}
\emph{yadi bhavān yathāśakti yatate, tarhi alam.}
\end{quote}

\begin{quote}
If you try as hard as you can, that is enough.
\end{quote}

\begin{cousinweb}
\textbf{\sk{यथाशक्ति}} (\emph{yathāśakti}) is built from \emph{yathā}, \textquotedblleft{}as\textquotedblright{}, and \emph{śakti}, \textquotedblleft{}power, ability\textquotedblright{}; MW gives \emph{yathā-śakti} \textquotedblleft{}according to power or ability, to the utmost of one's power\textquotedblright{} (MW).
\end{cousinweb}

\begin{grammarlens}[title={What you've built}]
One more word for saying what happens if.
\end{grammarlens}

\subsection*{Your turn}
\begin{itemize}
\raggedright
  \item \emph{Say it:} \emph{yathāśakti}
  \item \emph{Say it:} \emph{yathāśakti}, once more
  \item \emph{Say it:} \emph{yadi bhavān yathāśakti yatate, tarhi alam.} — If you try as hard as you can, that is enough.
  \item \emph{Recall:} say the Sanskrit for perseverance, then the Sanskrit for constancy, then say \emph{yathāśakti} again
\end{itemize}

\begin{tcolorbox}[breakable,colback=teal!4,colframe=teal!35!black,title={Before you move on}]
What does \sk{यथाशक्ति} mean? (\textquotedblleft{}As far as one can\textquotedblright{}.) Say it once more.
\end{tcolorbox}
\section[\texorpdfstring{\sk{अनुतिष्ठति} (anutiṣṭhati)}{अनुतिष्ठति (anutiṣṭhati)}]{\sk{अनुतिष्ठति} (anutiṣṭhati) — carries out}

\label{lesson:SA-C384-anutisthati}



\begin{quote}
Before the new one: say the Sanskrit for constancy, then the Sanskrit for as far as one can.
\end{quote}

\subsection*{You'll want to know: \sk{अनुतिष्ठति}}
\textbf{\sk{अनुतिष्ठति}} — \emph{anutiṣṭhati} — \textquotedblleft{}carries out\textquotedblright{}.

\begin{quote}
\sk{यदि} \sk{सः} \sk{योजनाम्} \sk{अनुतिष्ठति}, \sk{तर्हि} \sk{लाभः} \sk{भविष्यति।}
\end{quote}

\begin{quote}
\emph{yadi saḥ yojanām anutiṣṭhati, tarhi lābhaḥ bhaviṣyati.}
\end{quote}

\begin{quote}
If he carries out the plan, there will be a gain.
\end{quote}

Past (he …): \sk{सः} \sk{योजनाम्} \sk{अनुष्ठितवान्}, \emph{saḥ yojanām anuṣṭhitavān}.

\begin{cousinweb}
\textbf{\sk{अनुतिष्ठति}} (\emph{anutiṣṭhati}) is built on \emph{anu}, \textquotedblleft{}after\textquotedblright{}, and the root √sthā, \textquotedblleft{}to stand\textquotedblright{}; MW gives \emph{anu-ṣṭhā} \textquotedblleft{}to carry out, to perform, to practise\textquotedblright{} and \emph{anu-ṣṭhita} \textquotedblleft{}done, practised\textquotedblright{} (MW).
\end{cousinweb}

\begin{grammarlens}[title={What you've built}]
One more word for saying what happens if.
\end{grammarlens}

\subsection*{Your turn}
\begin{itemize}
\raggedright
  \item \emph{Say it:} \emph{anutiṣṭhati}
  \item \emph{Say it:} \emph{anutiṣṭhati}, once more
  \item \emph{Say it:} \emph{yadi saḥ yojanām anutiṣṭhati, tarhi lābhaḥ bhaviṣyati.} — If he carries out the plan, there will be a gain.
  \item \emph{Recall:} say the Sanskrit for constancy, then the Sanskrit for as far as one can, then say \emph{anutiṣṭhati} again
\end{itemize}

\begin{tcolorbox}[breakable,colback=teal!4,colframe=teal!35!black,title={Before you move on}]
What does \sk{अनुतिष्ठति} mean? (\textquotedblleft{}Carries out\textquotedblright{}.) Say it once more.
\end{tcolorbox}
\section[\texorpdfstring{\sk{प्रोत्साहयति} (protsāhayati)}{प्रोत्साहयति (protsāhayati)}]{\sk{प्रोत्साहयति} (protsāhayati) — encourages}

\label{lesson:SA-C384-protsahayati}



\begin{quote}
Before the new one: say the Sanskrit for as far as one can, then the Sanskrit for carries out.
\end{quote}

\subsection*{You'll want to know: \sk{प्रोत्साहयति}}
\textbf{\sk{प्रोत्साहयति}} — \emph{protsāhayati} — \textquotedblleft{}encourages\textquotedblright{}.

\begin{quote}
\sk{यदि} \sk{माता} \sk{प्रोत्साहयति}, \sk{तर्हि} \sk{बालः} \sk{उत्साहेन} \sk{पठति।}
\end{quote}

\begin{quote}
\emph{yadi mātā protsāhayati, tarhi bālaḥ utsāhena paṭhati.}
\end{quote}

\begin{quote}
If his mother encourages him, the child studies with zeal.
\end{quote}

Past (he …): \sk{सः} \sk{छात्रान्} \sk{प्रोत्साहितवान्}, \emph{saḥ chātrān protsāhitavān}.

\begin{cousinweb}
\textbf{\sk{प्रोत्साहयति}} (\emph{protsāhayati}) is built on \emph{pra}, \emph{ut} and the root √sah, \textquotedblleft{}to be able, to endure\textquotedblright{}, in its causative form; MW gives \emph{prot-sāhita} \textquotedblleft{}incited, stimulated, encouraged\textquotedblright{} (MW).
\end{cousinweb}

\begin{grammarlens}[title={What you've built}]
One more word for saying what happens if.
\end{grammarlens}

\subsection*{Your turn}
\begin{itemize}
\raggedright
  \item \emph{Say it:} \emph{protsāhayati}
  \item \emph{Say it:} \emph{protsāhayati}, once more
  \item \emph{Say it:} \emph{yadi mātā protsāhayati, tarhi bālaḥ utsāhena paṭhati.} — If his mother encourages him, the child studies with zeal.
  \item \emph{Recall:} say the Sanskrit for as far as one can, then the Sanskrit for carries out, then say \emph{protsāhayati} again
\end{itemize}

\begin{tcolorbox}[breakable,colback=teal!4,colframe=teal!35!black,title={Before you move on}]
What does \sk{प्रोत्साहयति} mean? (\textquotedblleft{}Encourages\textquotedblright{}.) Say it once more.
\end{tcolorbox}
